\documentclass[11pt,a4paper]{article}
\usepackage[utf8]{inputenc}
\usepackage[vietnamese]{babel}
\usepackage{amsmath,amssymb}
\usepackage{geometry}
\geometry{a4paper, total={180mm,260mm}, left=15mm, right=15mm, top=18mm, bottom=18mm}

\usepackage{tikz}
\usetikzlibrary{shapes.geometric, arrows.meta, positioning, calc}
\tikzset{
  arr/.style={-{Stealth[scale=0.8]}},
  every node/.append style={align=center}
}
\usepackage{tcolorbox}
\tcbuselibrary{skins,breakable}
\usepackage{fontawesome}
\usepackage{fancyhdr}
\usepackage{enumitem}

\pagestyle{fancy}
\fancyhf{}
\renewcommand{\headrulewidth}{0.5pt}
\renewcommand{\footrulewidth}{0.5pt}
\fancyhead[L]{\small\bfseries\color{violet!85!black}\faBook\ SỔ TAY KIẾN THỨC TOÁN 2}
\fancyhead[R]{\small\bfseries\color{violet!85!black}Lớp Toán Thầy Lĩnh -- loptoanthaylinh.vn}
\fancyfoot[C]{\small\color{gray}-- Sổ tay kiến thức: Bài 16. Lít (l) --}

\setlength{\parindent}{0pt}
\setlength{\parskip}{0pt}

\begin{document}

% --- BANNER BÀI HỌC ---
\begin{tcolorbox}[
  enhanced,
  colback=violet!8!white,
  colframe=violet!70!black,
  arc=3mm, boxrule=1.2pt,
  left=3mm, right=3mm, top=2mm, bottom=2mm,
  halign=center
]
  {\LARGE\bfseries\color{violet!90!black}\faBookmark\ BÀI 16: LÍT (L)}\\[1mm]
  {\itshape\footnotesize\color{black!75} Đơn vị đo dung tích lít, ca đong và kĩ năng tính toán dung tích chất lỏng trên 1 trang A4}
\end{tcolorbox}

\vspace{-2mm}

\noindent
\begin{minipage}[t]{0.49\textwidth}\vspace{0pt}
\begin{tcolorbox}[
  enhanced,
  colback=violet!4!white,
  colframe=violet!75!black,
  arc=2mm, boxrule=0.8pt,
  left=2mm, right=2mm, top=1.8mm, bottom=1.8mm,
  title={\footnotesize\bfseries\color{white}\faKey\ 1. Đơn vị đo dung tích lít}
]
{\footnotesize
\begin{itemize}\setlength{\itemsep}{1pt}\setlength{\leftskip}{-2.5mm}
  \item Để đo sức chứa (dung tích) của một vật (như chai nước, can dầu, bình nước, xô,...), người ta dùng đơn vị \textbf{lít}.
  \item Lít viết tắt là chữ \textbf{l} (thường in nghiêng $l$).
  \item \textbf{Dụng cụ đo:} Các ca đong $1\ l, 2\ l, 5\ l$; chai tiêu chuẩn $1\ l$,...
  \item \textbf{Thực tế:} Một chai nước ngọt lớn thường chứa $1\ l$ hoặc $1{,}5\ l$; một bình nước lọc trường học chứa $20\ l$.
\end{itemize}
}
\end{tcolorbox}
\end{minipage}
\hfill
\begin{minipage}[t]{0.49\textwidth}\vspace{0pt}
\begin{tcolorbox}[
  enhanced,
  colback=teal!4!white,
  colframe=teal!75!black,
  arc=2mm, boxrule=0.8pt,
  left=2mm, right=2mm, top=1.8mm, bottom=1.8mm,
  title={\footnotesize\bfseries\color{white}\faSliders\ 2. Phép tính với lít}
]
{\footnotesize
\begin{itemize}\setlength{\itemsep}{1pt}\setlength{\leftskip}{-2.5mm}
  \item Ta thực hiện cộng, trừ các số đo như các số tự nhiên rồi \textbf{viết thêm kí hiệu $l$ vào sau kết quả}.
  \item \textbf{Ví dụ:}
  \begin{itemize}\setlength{\itemsep}{0.5pt}\setlength{\leftskip}{-2.5mm}
    \item $9\ l + 5\ l = \mathbf{14\ l}$.
    \item $18\ l - 9\ l = \mathbf{9\ l}$.
    \item $15\ l + 5\ l - 8\ l = 20\ l - 8\ l = \mathbf{12\ l}$.
  \end{itemize}
  \item \textbf{Rót ra / Đổ thêm:}\\
  - Đổ thêm nước vào bình $\implies$ Làm tính \textbf{cộng} ($+$).\\
  - Rót bớt nước ra cốc $\implies$ Làm tính \textbf{trừ} ($-$).
\end{itemize}
}
\end{tcolorbox}
\end{minipage}

\vspace{-1.5mm}
\begin{tcolorbox}[
  enhanced,
  colback=orange!4!white,
  colframe=orange!80!black,
  arc=2mm, boxrule=0.8pt,
  left=2mm, right=2mm, top=1.5mm, bottom=1.5mm,
  title={\footnotesize\bfseries\color{white}\faSitemap\ 3. Mô hình trực quan: Rót nước và sức chứa ca đong}
]
\centering
\begin{tikzpicture}[x=1.5cm, y=1cm]
  % Can lớn
  \draw[thick, violet!80!black, fill=violet!15] (0,0) rectangle (1.2,1.8);
  \node[font=\tiny\bfseries] at (0.6, 0.9) {Can lớn\\$10\ l$};
  \draw[->, thick, red!85!black] (1.3, 0.9) -- node[above, font=\tiny\bfseries] {Rót sang} (2.3, 0.9);
  % Can nhỏ 1
  \draw[thick, teal!70!black, fill=teal!15] (2.5,0) rectangle (3.3,1.2);
  \node[font=\tiny\bfseries] at (2.9, 0.6) {Bình $1$\\$3\ l$};
  % Can nhỏ 2
  \draw[thick, teal!70!black, fill=teal!15] (3.6,0) rectangle (4.4,1.5);
  \node[font=\tiny\bfseries] at (4.0, 0.75) {Bình $2$\\$4\ l$};
  % Còn lại
  \node[rectangle, rounded corners=2pt, draw=orange!80!black, fill=orange!20, font=\scriptsize\bfseries] at (6, 0.9) {
    Còn lại trong can:\\$10 - (3 + 4) = \mathbf{3\ l}$
  };
\end{tikzpicture}
\end{tcolorbox}

\vspace{-1.5mm}
\noindent
\begin{minipage}[t]{0.49\textwidth}\vspace{0pt}
\begin{tcolorbox}[
  enhanced,
  colback=green!4!white,
  colframe=green!60!black,
  arc=2mm, boxrule=0.8pt,
  left=2mm, right=2mm, top=1.8mm, bottom=1.8mm,
  title={\footnotesize\bfseries\color{white}\faCheckCircle\ 4. Ví dụ thực hành chuẩn mực}
]
{\footnotesize
\textbf{Bài toán mẫu:}\\
Một thùng chứa $35\ l$ dầu, người ta đã bán đi $18\ l$ dầu. Hỏi trong thùng còn lại bao nhiêu lít dầu?\\
\textit{Bài giải:}\\
Trong thùng còn lại số lít dầu là:
$$35 - 18 = 17\text{ (}l\text{)}$$
Đáp số: $17\ l$ dầu.
}
\end{tcolorbox}
\end{minipage}
\hfill
\begin{minipage}[t]{0.49\textwidth}\vspace{0pt}
\begin{tcolorbox}[
  enhanced,
  colback=red!4!white,
  colframe=red!75!black,
  arc=2mm, boxrule=0.8pt,
  left=2mm, right=2mm, top=1.8mm, bottom=1.8mm,
  title={\footnotesize\bfseries\color{white}\faExclamationTriangle\ 5. Mẹo tránh bẫy trắc nghiệm}
]
{\footnotesize
\begin{itemize}\setlength{\itemsep}{1pt}\setlength{\leftskip}{-2.5mm}
  \item \textbf{Phân biệt đơn vị Khối lượng và Dung tích:}\\
  - Cân nặng (gạo, thịt, người): Đơn vị là \textbf{kg}.\\
  - Thể tích chất lỏng (nước, sữa, xăng): Đơn vị là \textbf{lít ($l$)}.
  \item Không cộng số đo khác đơn vị ($5\text{ kg} + 3\ l$ là vô nghĩa!).
  \item Luôn ghi kèm kí hiệu $l$ ở phép tính và đáp số.
\end{itemize}
}
\end{tcolorbox}
\end{minipage}


\end{document}
